\documentclass[11pt]{article}
\usepackage[margin=1in]{geometry}
\usepackage{amsmath,amssymb,amsthm,mathtools,longtable,array,hyperref}
\allowdisplaybreaks[2]
\setlength{\emergencystretch}{3em}
\newtheorem{theorem}{Theorem}[section]
\newtheorem{lemma}[theorem]{Lemma}
\newtheorem{proposition}[theorem]{Proposition}
\theoremstyle{remark}\newtheorem{remark}[theorem]{Remark}
\newcommand{\dd}{\,\mathrm d}
\newcommand{\E}{\mathbb E}
\newcommand{\KK}{\mathcal K}
\newcommand{\BB}{\mathcal B}
\newcommand{\Gp}{G_+}
\newcommand{\eps}{\varepsilon}
\title{Weak-passage assembly: a two-boundary theorem and the remaining response barrier}
\author{Analytical increment AN03 weak passage theorem v1}
\date{Version 1}
\begin{document}\maketitle
\begin{abstract}
The proved stopped implications strictly improve every proposed
energy, residual, field, core and outside guard. There is, however,
one unmatched domain hypothesis: the cited implications assume
$0\le c\le c_b$, whereas the proposed initialized first exit
addresses only the upper response boundary. We prove the resulting
two-boundary statement: starting at strong half mass, all the
assembled estimates hold until the first hit of either $0$ or $c_b$,
and one of these boundaries is reached within $20\log(1/q)$.
The sources do not exclude the lower hit. Thus the requested
unconditional passage to $c_b$ is not claimed. We isolate the exact
missing signed inequality and record the correct clock accuracy;
the denominator $Q(1-c_b)$ supplies a duration upper bound, not
the proposed asymptotic equality with an additive $O(1)$ numerator.
\end{abstract}

\section{Setting and scope of the result}
All new labels have prefix \texttt{an03wptv1:}. Source names are
those of roadmap v1.16; their exact paths and blob identifiers are
in the companion review. We use the initialized original balanced
flow, not a leading-system trajectory:
\[
 U(0,b)=W(0,b)=q^5(\cos b,\sin b),\qquad
 \dd\lambda_0=\dd b/(2\pi),\qquad
 P=\tfrac12N(e_1,q^2I)+\tfrac12N(\rho e_2,q^2I).
\]
Put $\lambda=\rho^2\in[169/400,49/100]$, $Q=\lambda+q^2$,
$\ell=\log(1/q)$, $k=10(1-\lambda)$ and $t_0=t_h$.
The response is the signed regression coefficient
\begin{equation}
 c(t)=Q^{-1}\E_{P_2}[f_2(X,t)X_2],\qquad
 I(t)=Q\int_{t_0}^t(1-c(s))\dd s.
 \label{an03wptv1:response}
\end{equation}
It is not a nonnegative mass by definition.

Use the fixed partition $\KK,\Gp,\Sigma,\BB$, with
$\mathcal D=\varnothing$, $G=\Gp\cup\Sigma$ and
$H=H_{\mathcal C_q}$ for the original $\eta=1/8$ clock cohort.
The core is the TAILREP/S2CP fixed-label set in the right-half
ancestry, selected at $t_0$ by
\begin{equation}
 \xi\le d_c=c_*q^{\chi+kP_\rho/2},\qquad
 \chi=1/20,\quad P_\rho=\Phi(\rho a),\quad a=\rho K(\rho a).
 \label{an03wptv1:core}
\end{equation}
Here $\xi$ is joint scaled displacement from MOMC's actual reference.
The complement of this core in the right-half ancestry is
$\mathcal C$. The entire outside set is exactly
\begin{equation}
 \BB=\mathcal C\cup\mathcal Q_{3,\mathrm{int}}^0
              \cup\mathcal S_\pi\cup\mathcal S_{3\pi/2},
 \label{an03wptv1:outsidepartition}
\end{equation}
up to null boundary labels, assigned once. No time-dependent
membership or new explicit cohort is introduced.

We retain the source energy notation $E_S=H_S+A_S$,
$\mu_S=E_S+J_S$, and write
\[
 Y=\sqrt{J_G}/q,\quad r=\sup_\theta|R_{11}|,\quad
 e=|(x_c,y_c)-(a,a)|_2,\quad
 r_{\KK}=\mathop{\rm ess\,sup}_{b\in\KK}
       (|x_b-x_c|+|y_b-y_c|).
\]
The reference is never reset. Let
\begin{equation}
 t_-:=\inf\{t\ge t_0:c(t)=0\},\qquad
 t_+:=\inf\{t\ge t_0:c(t)=c_b\},\qquad
 \tau:=\min(t_-,t_+),
 \label{an03wptv1:stops}
\end{equation}
with the infimum of an empty set equal to infinity.

\begin{theorem}[Initialized assembly up to the first response boundary]
\label{an03wptv1:twoboundary}
There are fixed constants, chosen in Section 2's order, and a
small fixed $c_b>0$ such that, for all sufficiently small $q$
uniformly on the ratio box,
\[
 0<c(t_0)<c_b,\qquad t_0<\tau<t_0+20\ell.
\]
On the entire closed interval $[t_0,\tau]$ the following hold:
\begin{align}
 \|e_1\|_{P_1}+\|e_2\|_{P_2}&\le C,&
 \|e_2\|_{P_1}+\|e_1\|_{P_2}&\le M_{\rm res}q/2,
 \label{an03wptv1:residuals}\\
 J_G&\le q^2K_Y^2/4,&
 \int_{t_0}^\tau r&\le C_r\int_{t_0}^\tau|1-\kappa_1|+o(1),
 \label{an03wptv1:transverse}\\
 \int_{t_0}^\tau|1-\kappa_1|&\le\log2+o(1),&
 \mu_{\KK}(t)&\le e^{d_q},\quad d_q=o(1),
 \label{an03wptv1:learning}\\
 H\le H_{\Gp}&\le\tfrac98c_b+o(1),&
 E_{\Gp}&\le\tfrac54c_b+o(1),
 \label{an03wptv1:energies}\\
 E_\Sigma&\le C_\Sigma(q^5+q^{\Gamma_\nu}H),&
 \Gamma_\nu&\ge289169/924000>3/10,
 \label{an03wptv1:secondary}\\
 r_{\KK}(t)&=O(q^{1/20}),& e(t)&\le3e_*/4.
 \label{an03wptv1:coreconclusion}
\end{align}
The core remains in the fixed coordinate chart
$|\theta|,|\psi|<Rq$, $R=79/51$. Uniformly in time,
\begin{align}
 \mu_{\mathcal C}/q&\le Cq^{\Delta_\nu},
       &\Delta_\nu&>3/25,\nonumber\\
 \mu_{\BB}/q&\le C\{q^{\Delta_\nu}+q^9
      +s_qq^{411/800}+(1+s_q)q^{769/4000}
      +q^{10}\sqrt\ell\}=o(q^{1/10}),
 \label{an03wptv1:outside}
\end{align}
where $s_q=O(\sqrt{\log\ell})$. The corresponding time integrals
are bounded by $20\ell$ times these right sides. With the canonical
CFA c-form remainder and forcing,
\begin{align}
 \int_{t_0}^\tau\eps_F
 &\le C_FC_{L_c}C_{\rm RCET}C_0E_g+o(1)
        <\overline F_F/2,
 \label{an03wptv1:field}\\
 \int_{t_0}^\tau f_*
 &\le\frac{10e}{9}\frac{E_g}{Q(1-c_b)}
  +\frac{C_V(E_g+K_Y\sqrt{E_g})}{Q(1-c_b)}
  +C_FC_{L_c}C_{\rm RCET}C_0E_g+o(1)
  \le\frac{e_*}{4C_{\rm pass}}.
 \label{an03wptv1:forcing}
\end{align}
Here $E_g=2c_b$, $f_*=c+|V|+K_w+\eps_F$, and $T\equiv0$.
The clock is positive and satisfies, for all subintervals,
\begin{equation}
 \log\frac{H(t)}{H(s)}=I(t)-I(s)+\delta(s,t),\qquad
 |\delta(s,t)|\le D_{\rm clk}
 =O(q^{1/8}+q^2\ell)=o(1).
 \label{an03wptv1:clock}
\end{equation}
The sign-good group has RCET's bounded defect
$O(C_0+q^2\ell)$ and fixed-factor relative-amplitude transport.

The theorem does not identify which response boundary is hit.
In particular it does not assert $\tau=t_b=t_+$. A nonnegative-response
barrier, or the stronger exclusion $t_->t_+$, remains to be supplied
as specified in Section 6.
\end{theorem}

\section{Constants and the canonical field guard}
Fix at the start
\begin{equation}
 \bar c=1/4,\quad L_T=20,\quad b_B(q)=q^{1/10},\quad
 R=79/51,\quad M_K=2,\quad B_s^0=\log2+1,
 \quad L_c=R+1.
 \label{an03wptv1:firstconstants}
\end{equation}
Choose a fixed $c_*>0$ and $0<\nu\le1/100$ as in S2CP.
The core norm convention needs one explicit adjustment: LCRC writes
$|\theta|+|\psi|\le R_{\rm LCRC}q$, whereas MOMC/CFA use
coordinate bounds $Rq$. Use $R_{\rm LCRC}=2R$ in LCRC's
constants. This is a fixed model radius chosen now; it does not
enlarge CFA's receiver square or create a residual-radius loop.

Use CFA's model-uniform $C_{\rm pass},g,e_*,K_{\rm pass}$ and set
$P_{\rm allow}=e_*/(4C_{\rm pass})$. Following CFA's uniform
construction, use the conservative energy $2+2\bar c$, amplitude
$\bar h_*=(9/8)\bar c+1$ and lower rate
$\underline b=(169/400)(1-\bar c)$ in LCRC's formulas for
$E_*,K_E,K_{1/2},A_*,K_Y,M_{\rm res}$. In particular
\[
 B=C_rB_s^0+2.
\]
Enlarge $M_{\rm res}$ once, if necessary, so that
$M_{\rm res}>2C_{\rm pre}$ from Q3SB. All these constants are
independent of the subsequently selected $c_b,C_0$.

Choose any fixed positive coarse field allowance $\overline F_F$.
Choose $0<c_b\le\bar c$ small enough that the first two terms
of \eqref{an03wptv1:forcing}, with $E_g=2c_b$ and $Q\ge169/400$,
sum to at most $P_{\rm allow}/4$. Fix S2CP's $C_D,A_0,C_m$
from these constants and $\overline F_F$. Choose $C_0$ within
LCRC/RCET's admissibility and CFA's fixed-chart boundary margin,
with
\begin{equation}
 C_FC_{L_c}C_{\rm RCET}C_0E_g
 \le\min(\overline F_F/4,P_{\rm allow}/4).
 \label{an03wptv1:C0choice}
\end{equation}
Only then fix $K(C_0)$, SEP's cap and entry constants, and
$C_\Sigma\ge\max(1,C_{\rm out},C_{\rm ent}+1)$ as in S2CP.
Choose $q$ last. This is CFA's order, with duration and the outside
function made explicit at the first step. No positive-power error
coefficient is used to choose an earlier order-one constant.

The field guard is a well-defined scalar quantity, not an existence
assertion. For every state in the core chart, use CFA's exact
definitions
\[
 \widetilde\mu_s=\mu_s\min(1,\mu_{\KK}^{-1}),\quad
 V=V_O,\quad K_w=(K_O)_+,\quad T=0,
 \quad\eps_F(t)=\|F_{\rm orig}-F_c\|_*.
\]
Here $F_c$ is CFA \texttt{cform}, $\|\cdot\|_*$ controls value and
induced $\ell^1$ receiver derivative on the fixed square, and
$O=\Gp\cup\Sigma\cup\BB$. The remainder is defined by subtraction
before its budget is improved. The normalized donor measure is
positive and has mass at most one. Thus imposing
$\int\eps_F<\overline F_F$ supplies precisely the inherited
S2CP field hypothesis, without first assuming CFA's sharper bound.

\section{Strict entry, including the response}
\begin{lemma}[The initialized entry lies inside the joint domain]
\label{an03wptv1:entry}
For the above fixed constants and sufficiently small $q$, all guards
in Section 4 hold strictly at $t_0$. In addition,
\begin{equation}
 c(t_0)=\frac{a^2}{2\lambda}q^2+o(q^2)>0.
 \label{an03wptv1:centry}
\end{equation}
\end{lemma}
\begin{proof}
MOMC supplies the actual reference $e_0=o(1)$, right-half mass
$1/2$ and the raw right-half entry tail. The latter gives
$\mu_{\mathcal C}(t_0)/q\le Cq^{\Delta_\nu}$: the entry exponent
in S2CP's tail calculation is at least its later exponent
$\Delta_\nu>3/25$. Hence $\mu_{\KK}(t_0)=1/2-o(1)$.
The core definition gives $r_0\le d_c=o(1)$.
Q2EC/Q2EP give $J_G(t_0)\le C_Jq^3$,
$E_{\Gp}(t_0)=o(1)$, the primary signs and ratios, and
$H_0\asymp A_q^{-\lambda}q^k$, with $A_q$ bounded above and below.
SEP gives $E_\Sigma(t_0)\le C_{\rm ent}q^5<C_\Sigma q^5$.
Q3SB gives the cross-residual entry bound $C_{\rm pre}q$.
All cumulative-integral guards are zero at entry.

To check outside entry without using a post-entry guard, use the
pre-entry parts of the strip proofs. GESG's original Q3 interior
has mass $O(q^{10})$. PISB's outer and middle bounds, restricted
to $[0,t_0]$, use loss and Q3SB's proved pre-entry cross bound only.
PISB \texttt{entrymass} bounds an inner $\pi$ label by
\[
 C A_q^{-1}\{q^{c_m}+q^{2+(1-\lambda)c_m}+q^{4/\lambda-2}\}e^{o(1)},
 \quad c_m=4/(1-\lambda)+1/2.
\]
Each displayed exponent exceeds six; its label width is $qs_q$.
LSB \texttt{entrymass} and \texttt{density}, whose proofs end at
$t_0$ using only S1, give
\[
 \mu_{J_3}(t_0)\le C A_q^{-1}(q^{c_m+1}s_q+q^3)=O(q^3).
\]
Thus the whole Q3 entry mass is $O(q^3)$; in particular every
outside component is strictly below $q b_B(q)/8$ for small $q$.
This verifies the outside guard independently of its own assumed
post-entry version.

There is also a short signed entry check. MOMC's uniform shapes
at $t_0$ put every right-half donor at
$(\theta/q,\psi/q)=(a,a)+o(1)$, with mass $1/2$.
Under $P_2$, write $X=(qZ,V)$, $V=\rho+qZ_2$.
The same compact Gaussian moment used in MOMC's reduction gives
\[
 c_{\mathcal R}(t_0)
 =\frac{q^2}{Q}\int_{\mathcal R}
 y\,\E[V(Z+xV)_+]\dd\mu+o(q^2)
 =\frac{a\rho K(\rho a)}{2\lambda}q^2+o(q^2).
\]
Use $a=\rho K(\rho a)$. The whole Q2 entry mass is $O(q^3)$
by Q2EP, and we just checked the same bound for Q3.
Their absolute responses are $O(q^3)$ by LCRC's elementary
mass charge. This proves \eqref{an03wptv1:centry}, uniformly in
the ratio box; $a^2/\lambda$ has a positive model lower bound.
The upper response boundary is fixed, so $c(t_0)<c_b$ as well.
\end{proof}

\section{Joint stopped interval and hypothesis matching}
In addition to stopping at either response boundary and at
$t_0+20\ell$, stop at the first equality in
\begin{equation}
 \begin{gathered}
 Y<K_Y,\quad\int r<B,\quad E_{\Gp}<2c_b,\quad
 \|e_2\|_{P_1}+\|e_1\|_{P_2}<M_{\rm res}q,\\
 E_\Sigma<2C_\Sigma(q^5+q^{\Gamma_\nu}H),\qquad
 \int\eps_F<\overline F_F,\\
 e<e_*,\quad r_{\KK}<1,\quad
 |\theta|,|\psi|<Rq\text{ on }\KK,\quad\mu_{\KK}<2,\\
 \sup_{t_0\le s\le t}\mu_{\BB}(s)/q<b_B(q).
 \end{gathered}
 \label{an03wptv1:guards}
\end{equation}
On each such initial interval $0<c<c_b$. Its closed endpoint
has $0\le c\le c_b$. The sources apply with that endpoint;
none needs the future time $t_b$ to be known in advance.

\begin{proposition}[Every non-response guard improves]
\label{an03wptv1:improvements}
On the joint stopped interval, the energy, residual, field, core
and outside guards in \eqref{an03wptv1:guards} have strict slack,
uniformly in its endpoint. They cannot cause the first exit.
\end{proposition}
\begin{proof}
We match hypotheses; the source arguments themselves are not
reproved. The initially assumed core geometry supplies both
S2CP's and LCRC's core state inputs. S2CP Proposition Bs therefore
gives
\[
 \int|1-\kappa_1|\le\log2+o(1),\qquad
 (\log\mu_{\KK})_+(t)+\int(\mu_{\KK}-1)_+\le d_q=o(1).
\]
This proposition uses the current core geometry and instantaneous
outside guard, not future core retention. In particular it supplies
LCRC's learning input $B_s^0$ and improves $\mu_{\KK}<2$.

LCRC \texttt{localclock} uses only entry, response domain, loss,
cross residual, duration and $\int r<B$. It gives the primary
signs, $E_{\Gp}\le(10/9)H_{\Gp}$, its bounded backward defect,
and the separate $o(1)$ clock defect. It uses neither a secondary
energy conclusion nor an endpoint cap. S2CP \texttt{gate} then
uses the temporary secondary guard, core and outside state guards
to bound the signed gate excess. S2CP Theorem S2 additionally
uses the coarse field guard and reference bound; SEP and MOMC
supply its entry profiles. Consequently it gives its distance,
gain, secondary energy and strong-complement mass conclusions.
Its output $C_\Sigma$ was chosen after the entry constants and
does not depend on the coefficient of the temporary secondary
guard. Thus the factor-two secondary guard improves strictly.
Its Corollary U gives $\mathcal U_\Sigma=o(1)$.

With those distance/gain conclusions, all LCRC retained secondary
inputs are supplied on the very same interval. The chart-group
condition and \texttt{chartmass} are vacuous because $\mathcal D$
is empty. The pointwise ledger cap yields
$H_{\Gp}\le(9/8)c_b+o(1)$ and
$E_{\Gp}\le(5/4)c_b+o(1)<(3/2)c_b<2c_b$.
LCRC's unchanged Volterra step gives $Y\le K_Y/2$ and its
instantaneous output estimate gives cross residual at most
$M_{\rm res}q/2$. Its integral residual estimate and Bs give
$\int r\le C_r\log2+o(1)<B-1$.
The conservative constants used before selecting $c_b$ suffice
for this smaller energy guard, exactly as checked in S2CP and CFA.

For outside mass, S2CP \texttt{compmass}, GESG's original-label
Q3 interior, PISB \texttt{main} and LSB \texttt{main} give
\eqref{an03wptv1:outside} under the current stopped hypotheses.
Their pre-entry cross input is Q3SB \texttt{crosspre}.
PISB and LSB require neither CFA's strip budget nor its passage
conclusion. For each of the four fixed components in
\eqref{an03wptv1:outsidepartition}, its power is strictly larger
than $1/10$. All coefficients are fixed before $q$; hence each
is at most $b_B(q)/8$, and their sum is at most $b_B(q)/2$.
Their integrated versions supply CFA's strip input with
$F_{\mathcal S}(q)=o(1)$, satisfying any prescribed fixed small
allowance for sufficiently small $q$.

CFA \texttt{assembly} now uses S2CP's energy, outside and overshoot
conclusions, the current core geometry and the paid strip input.
It improves the canonical field guard to
\[
 \int\eps_F\le C_FC_{L_c}C_{\rm RCET}C_0E_g+o(1)
       <\overline F_F/2.
\]
The coefficient of the primary term is
independent of the coarse field allowance and secondary guard
coefficient. Their influence in the other source constants
multiplies only vanishing scales. CFA \texttt{forcing} gives
\eqref{an03wptv1:forcing}, with its factor $1/Q$ retained.

Finally CFA \texttt{passage} applies at this stopped endpoint.
Its amplitude cap follows from LCRC's pointcap, its signed clock
defect from localclock/BULK, its initial $e_0,r_0,H_0$ from S1
and \eqref{an03wptv1:core}, and its forcing budget was just paid.
Its cutoff calculation keeps
\[
 r_0L_b\le Cq^{\chi+kP_\rho q^2/(2Q)}=O(q^{1/20}).
\]
Thus $e\le3e_*/4$ and $r_{\KK}=O(q^{1/20})<1/2$.
Since $a<28/51$, $e_*\le1/4$ and $r_{\KK}=o(1)$, each core
coordinate is at most $28/51+3/16+o(1)<79/51$ in absolute
value. The separate core chart guard therefore improves too.
No future retained state has been used to prove its own past
hypothesis: each source was invoked on the joint stopped interval.
\end{proof}

\paragraph{Continuity and continuation.}
For each fixed positive $q$, the foundational characteristic
regularity and finite-horizon mass bounds give continuous dependence
on time uniformly over the compact initial circle. The fixed-cohort
energies and response are continuous by domination. Gaussian
integration gives continuous residual norms and matrix entries;
the supremum over the compact receiver circle is continuous.
The core is a fixed entry sublevel set. Uniform-in-label continuity
also gives continuity of its essential-supremum radius and chart
coordinates. The canonical field and its first receiver derivative
are continuous on the fixed chart; the mass cap and positive-part
moment in its definition are continuous, even at their corners.
Thus $\eps_F$ is locally bounded and its time integral continuous.
The cumulative outside supremum is continuous as well. All guards
start strictly and all non-response first equalities are excluded
by Proposition~\ref{an03wptv1:improvements}. This justifies the
common continuation argument, including a simultaneous exit.

\section{Duration and the clock accuracy actually certified}
LCRC's ledger and sign-good estimate give on a stopped interval
\begin{equation}
 c(t)\ge\frac89H_{\Gp}(t)-\epsilon_q,\qquad
 \epsilon_q=C\{q^2+q b_B(q)+q^5+q^{\Gamma_\nu}E_g
       +q^2K_Y\sqrt{q^5+q^{\Gamma_\nu}E_g}\}+o(1),
 \label{an03wptv1:ledgerlower}
\end{equation}
where the additional $o(1)$ includes the exponentially small
primary tail. The constants are fixed. Combining with the clock,
\begin{equation}
 e^{-D_{\rm clk}}H_0e^{Q(1-c_b)(t-t_0)}
 \le H(t)\le H_{\Gp}(t)\le\frac98(c_b+\epsilon_q).
 \label{an03wptv1:durationineq}
\end{equation}
Therefore every endpoint inside this response strip satisfies
\begin{equation}
 t-t_0\le
 \frac{\log(9c_b/(8H_0))+D_{\rm clk}+1}{Q(1-c_b)}
 \quad\text{for sufficiently small }q.
 \label{an03wptv1:duration}
\end{equation}
This is the proposed duration argument with its correct scope:
it excludes remaining in $[0,c_b]$ for too long, not exit at zero.

Since $H_0\asymp A_q^{-\lambda}q^k$, $k\le231/40$ and
$Q(1-c_b)\ge507/1600$,
\[
 20(507/1600)-231/40=9/16>0.
\]
Thus the right side of \eqref{an03wptv1:duration} is less than
$20\ell$ for small $q$, independently of all later guards.
The common stop cannot be the duration boundary. The strict
improvements and entry lemma prove
Theorem~\ref{an03wptv1:twoboundary}.

There is a sharper, useful clock identity. At any certified endpoint,
\begin{equation}
 t-t_0=\frac{\log(H(t)/H_0)-\delta(t_0,t)}Q
                         +\int_{t_0}^t c(s)\dd s,
 \quad |\delta(t_0,t)|\le D_{\rm clk}.
 \label{an03wptv1:exacttime}
\end{equation}
S2CP's gate budget gives
\begin{equation}
 0\le\int c\le\frac{10e}{9}\frac{E_g}{Q(1-c_b)}+o(1).
 \label{an03wptv1:cint}
\end{equation}
In particular, using the cap on $H$,
\begin{equation}
 t-t_0\le\frac{\log(c_b/H_0)+O_{
m fixed}(1)}Q.
 \label{an03wptv1:sharperupper}
\end{equation}
The requested equality
$t_b-t_0=[\log(c_b/H_0)+O(1)]/[Q(1-c_b)]$ is not the
accuracy of these sources. For any fixed $c_b>0$ it is incompatible
with \eqref{an03wptv1:sharperupper} if a certified $t_b$ exists:
the difference of the two leading terms is
\[
 \frac{c_b}{Q(1-c_b)}\log(c_b/H_0)
 =\frac{c_b k}{Q(1-c_b)}\ell+O_{
m fixed}(1),
\]
which diverges. Formula \eqref{an03wptv1:duration} is a valid
upper bound. Formula \eqref{an03wptv1:exacttime} retains the
$o(1)$ clock defect and the exact terminal clock amplitude.
An additional terminal comparison $H(t_b)\asymp c_b$, if established,
would give $[\log(c_b/H_0)+O(1)]/Q$, not the proposed denominator.
No such comparison is silently used here.

\section{Exact stopping point: the lower response boundary}
The hypotheses $0\le c\le c_b$ occur explicitly in LCRC Section 1,
RCET Theorem 2.1, BULK Theorem 4.1, S2CP's guard list, CFA's guard
list and CLOCK's cfield/passage statements. They are not conclusions
of those statements. The proposed guard list and the roadmap's
strict-improvement list omit an argument excluding exit through
the lower boundary. Continuity alone cannot exclude it.

The exact missing domain conclusion is
\begin{equation}
 \boxed{\quad c(t)\ge0\quad
   (t_0\le t\le\min\{t_+,t_0+20\ell\}).\quad}
 \label{an03wptv1:gap}
\end{equation}
This is needed for the initialized flow, with the other guards
improved as above. Touching zero and returning to positive response
would be allowed; strict positivity is not part of the required
contract. The stronger statement $t_->t_+$ would suffice to remove
the two-boundary alternative in our continuation argument.
A sufficient signed ledger estimate at a putative first
lower hit $s$ would be
\begin{equation}
 c_{\Sigma}(s)+c_{\KK}(s)+c_{\BB}(s)
 >-\frac89 H_{\Gp}(s)
       +C\mu_{\Gp}(s)e^{-\rho^2/(4q^2)}.
 \label{an03wptv1:signedneeded}
\end{equation}
Together with LCRC's one-sided primary tail this would imply
$c(s)>0$, a contradiction. A different proved invariant-region
argument for the exact response could also supply
\eqref{an03wptv1:gap}. Neither is among the matched sources.

The proposed absolute-charge route gives only
\eqref{an03wptv1:ledgerlower}. At $c(s)=0$ it implies
\[
 H_{\Gp}(s)\le\tfrac98\epsilon_q,
\]
which is compatible with a positive but small clock. To turn this
route into a lower-boundary contradiction one would need
\begin{equation}
 \frac89H_{\Gp}(s)>
 |c_\Sigma(s)|+|c_{\KK}(s)|+|c_{\BB}(s)|
       +C\mu_{\Gp}(s)e^{-\rho^2/(4q^2)}.
 \label{an03wptv1:absolutefailure}
\end{equation}
This is not a consequence of a vanishing absolute error. Indeed
at entry Q2EP gives $H_{\Gp}(t_0)\le\mu_{Q2}(t_0)=O(q^3)$,
whereas \eqref{an03wptv1:centry}'s core calculation also gives
$c_{\KK}(t_0)\asymp q^2$. Thus the uniform absolute-domination
version of \eqref{an03wptv1:absolutefailure} already fails at entry,
even though the actual response there is strictly positive.
The sign of the core contribution, lost in that absolute charge,
matters. Later the available outside bounds are also absolute and
are not $o(q^2)$ bounds on its response.

Q3SB \texttt{gatepre} bounds $|c|+N_-$ before $t_0$; it does not
prove nonnegativity. Its signed-excess calculation explicitly allows
slightly negative pre-entry response. Loss monotonicity and the
positive radial weights do not turn the signed regression
\eqref{an03wptv1:response} into a nonnegative mass. None of these
facts proves the missing later barrier.

We stop at \eqref{an03wptv1:gap}. This is an unmatched hypothesis,
not a counterexample to response positivity in the initialized flow.
No extension of CLOCK, BULK or the c-form inequalities to negative
response is asserted. Such an extension would be a different task.

\section{Handoff and boundary of the claim}
The certified handoff is at $\tau$, with the unresolved alternative
$c(\tau)=0$ or $c(\tau)=c_b$. If the latter alternative is proved
to occur, the following are available to S4 immediately:
\begin{itemize}
 \item Original residual bounds and $J_G=O(q^2)$, the primary
 sign/ratio transport with fixed $e^{O(C_0)}$ accuracy, and
 $\int r\le C_r\log2+o(1)$.
 \item The fixed clock seed $H_0\asymp A_q^{-\lambda}q^k$, the
 exact clock \eqref{an03wptv1:clock} with $o(1)$ defect, its
 amplitude cap, and the phase identity \eqref{an03wptv1:exacttime}.
 A fixed seed comparison is not an asymptotically exact seed coefficient.
 \item The same-field reference and fixed core, with
 $r_{\KK}=O(q^{1/20})$, $e\le3e_*/4$ and
 $\mu_{\KK}\le1+o(1)$; these are retention bounds, not an
 orbit-selection theorem.
 \item The secondary energy bound and all-endpoint capped gain;
 the complementary and outside masses, including their integrated
 $\mu/q$ budgets, at positive-power accuracy.
 \item The exact c-form representation, a small fixed integrated
 value-and-derivative remainder, and a small fixed CLOCK forcing
 budget. These budgets are not relabeled $o(1)$ when their
 $C_0E_g$ or $\sqrt{E_g}$ terms remain fixed.
\end{itemize}
As the sources stand, the condition selecting this upper-boundary
handoff has not been proved. No S4 weak return, matching beyond
$c_b$, or selection is begun, and no prior source is modified.
\end{document}
